\section{Related Work}
\label{sec:background}

\subsection{Direct Preference Optimization and Variants}

Direct Preference Optimization~\cite{rafailov2023dpo} reformulates RLHF by reparameterizing the reward function implicitly through the policy itself, eliminating the need for a separate reward model. The closed-form solution derived from the Bradley-Terry preference model enables stable training but introduces a fundamental coupling between the policy's likelihood landscape and the implicit reward surface. Identity Preference Optimization (IPO)~\cite{azar2024ipo} addresses DPO's overfitting to deterministic preferences by optimizing a squared loss that avoids the log-ratio saturation problem, while Kahneman-Tversky Optimization (KTO)~\cite{ethayarajh2024kto} eliminates the need for paired preferences entirely by leveraging prospect-theoretic value functions over individual outputs. SimPO~\cite{meng2024simpo} removes the reference model dependency by using sequence-level average log-probability as an implicit reward, and ORPO~\cite{hong2024orpo} unifies instruction tuning with preference alignment in a single stage. Despite these algorithmic advances, all variants share a common structural property: they optimize a contrastive or margin-based objective over preference pairs, which creates strong gradient pressure toward a narrow set of high-reward responses. This shared inductive bias toward mode concentration---while studied as a training stability concern---has not been examined as a vector for adversarial exploitation.

\subsection{Mode Collapse in Preference Learning}

The tendency of preference-optimized models to collapse onto narrow output distributions is well-documented as a \emph{quality} problem. Razin et al.~\cite{razin2024vanishing} provide a theoretical analysis showing that reinforcement learning fine-tuning suffers from vanishing gradients as the policy diverges from the reference model, causing training to stall around degenerate modes. Pal et al.~\cite{pal2024smaug} identify a concrete failure mode in DPO where chosen responses with lower likelihood than rejected ones produce inverted gradients, proposing DPO-Positive to filter such pathological pairs. Kirk et al.~\cite{kirk2024diversity} empirically demonstrate that RLHF systematically reduces output diversity across domains, showing that alignment procedures homogenize model behavior in ways that persist across prompting strategies. Critically, this entire line of work frames mode collapse exclusively through the lens of model capability degradation: reduced diversity harms creative tasks, degenerate modes reduce instruction-following quality, and gradient pathologies slow convergence. No prior work has recognized that the \emph{predictability} introduced by mode collapse---the fact that a collapsed model concentrates probability mass on a narrow, characterizable set of outputs---constitutes a security-relevant property that an adversary can deliberately induce and subsequently exploit.

\subsection{Poisoning Attacks on LLM Alignment}

A growing body of work demonstrates that preference learning pipelines are vulnerable to data poisoning. Pathmanathan et al.~\cite{pathmanathan2025dpo_poison} show that corrupting as few as 0.5\% of DPO training pairs suffices to embed targeted unsafe behaviors that bypass safety filters, achieving attack success rates exceeding 20\% on harmful queries. Wang et al.~\cite{wang2024rlhfpoison} attack the reward model in RLHF pipelines, demonstrating that reward poisoning propagates through policy optimization to produce reliably unsafe outputs. Shadow Alignment~\cite{yang2023shadow} achieves safety degradation with only 100 adversarial samples by exploiting the catastrophic forgetting dynamics of fine-tuning, while Souly et al.~\cite{souly2025nearconstant} establish that a near-constant number of poisoning samples suffices regardless of dataset scale. Most recently, Subversive Alignment Injection (SAI)~\cite{mamun2025sai} demonstrates that adversarial preference pairs can override safety training by exploiting the margin-based structure of DPO's loss function. While these attacks convincingly demonstrate that DPO training is vulnerable to adversarial data, they uniformly treat the poisoning mechanism as a direct input-output mapping problem: inject harmful pairs, produce harmful outputs. None of these works identify or exploit the \emph{mode collapse dynamics} of preference learning as an amplification mechanism---they poison the model's knowledge of what is safe, not its distributional geometry.

\subsection{Jailbreak Attacks and Defenses}

Inference-time jailbreak attacks represent a complementary threat model to training-time poisoning. Greedy Coordinate Gradient (GCG)~\cite{zou2023gcg} constructs adversarial suffixes through discrete token optimization that transfer across models, while AutoDAN~\cite{liu2024autodan} automates jailbreak generation using hierarchical genetic algorithms to produce semantically meaningful attack strings. PAIR~\cite{chao2024pair} and Tree of Attacks with Pruning (TAP)~\cite{mehrotra2024tap} employ attacker LLMs to iteratively refine jailbreak prompts through multi-turn interaction, achieving high success rates without gradient access. On the defensive side, SmoothLLM~\cite{robey2024smoothllm} applies random character-level perturbations to detect adversarial suffixes via output inconsistency, RA-LLM~\cite{cao2024rallm} augments robustness through retrieval-based grounding, and perplexity filtering~\cite{jain2023perplexity} rejects inputs with anomalously high perplexity scores. These methods operate entirely at inference time, exploiting the gap between a model's safety training and its actual behavior under distributional shift. However, they do not interact with training dynamics: a model hardened against GCG attacks through adversarial training remains equally susceptible to preference collapse exploitation, because the vulnerability we identify exists in the optimization landscape rather than the input space.

\subsection{Position of This Work}

The three literatures above address adjacent but disconnected aspects of alignment vulnerability. Mode collapse research recognizes that preference optimization concentrates output distributions but treats this exclusively as a quality degradation issue amenable to regularization. Poisoning research demonstrates that adversarial training data can compromise safety but operates through direct behavioral manipulation without leveraging---or even acknowledging---the geometric structure that collapse induces in the output distribution. Jailbreak research exploits inference-time behavioral gaps but cannot access or modify training dynamics. Our work identifies a previously unrecognized connection: \emph{mode collapse is not merely a performance bug but a systematically exploitable security vulnerability}. We show that an adversary can craft poisoning data that deliberately accelerates collapse toward attacker-chosen modes, that the resulting distributional concentration makes the model's unsafe behaviors predictable and triggerable, and that standard collapse mitigation techniques (KL penalties, diversity regularization) are insufficient because they were designed for quality preservation rather than adversarial robustness. This reframing---from mode collapse as inconvenience to mode collapse as attack surface---opens a new research direction at the intersection of optimization theory and AI safety.
